\section{Fourier inversion and the pointwise bound}
\label{sec:fourier}

At a real scale \(x\ge x_*\) of the fixed classical family,
Section~\ref{sec:integral} gives the signed mean bound
\eqref{eq:classical-averaged-excess} over \(H\).
Insert that bound into the Fourier calculation at this same \(x\).
Use the Artin data and character estimates of the seed \(i_0\), and
keep \(T,T_2,(a_t,b_t)_{t\in T},N\), the actual weights \(w_{u,j}\)
at the matched records \(i(u)\), and the centering in
\eqref{eq:centered-excess-definition}. The sampling estimates and
auxiliary character bounds apply on these data.

For the general fixed-target calculation, fix the target, precision
and family as in Sections~\ref{sec:family} and \ref{sec:assembly}.
The calculation uses the corresponding sampling and character
estimates and a bound on the signed mean over \(H_0\). In the
Dirichlet extension that mean bound follows from
\eqref{eq:integral-excess} under its lattice hypotheses.

\subsection{Sampling degrees and weights at the tied indices}
\label{sec:fourier-sampling}
\teamanchor{proof.fourier-sampling}

Write \(n=|H|\), \(f=2n\), \(N_T=|T|\), \(N_2=|T_2|\), and
\(a_E=xb\) with \(b=\exp(x^\varepsilon)\). In the chosen records the
other scales are
\[
 X'=\exp(x^{\varepsilon/3}),\qquad b'=\frac{a_E}{X'},\qquad
 w''=\frac{b'}{xn},\qquad \delta_1=x^{-\varepsilon/8}.
\]
The sampling integers \(a_t,b_t\) are read at \(i_0\) and carried by the
same numerical label \(t\in T\) to each \(i(u)\). The
\hyperlink{proof.selection-realized-samples}{admissibility inequality}
holds at every index over a sufficiently large scale of the selected family.
In particular the arithmetic base \(i_0\) itself
can supply the two arrays. The covering property then supplies an index
with these arrays at each prescribed translate.
\releasecite{Formalism/SharpRestriction/GlobalSchurBounds/WeightVectorSmallSharesEight.lean}{22}{Choice of the base arrays}
\releasecite{Formalism/SharpRestriction/FamCore.lean}{132}{Transport of both arrays}.
Their prescribed intervals are
\[
 a_E\le a_t\le a_E+\frac{a_E}{kx^{13}},\qquad
 b'\le b_t\le b'\left(1+\frac{2n\delta_1x^\varepsilon L_{\mathrm{a},G}}{k}\right),
 \qquad L_{\mathrm{a},G}\le x^{\varepsilon+\varepsilon/100}.
\]
Here \(L_{\mathrm{a},G}\) is the parameter in the upper endpoint of the
interval for \(b_t\). The degree window gives \(0<a_t\le2a_E\) and,
with \(y_t=a_t/b\), gives \(0\le y_t-x\le 1/(kx^{12})\le1\).
The representative \(i(0)\) may differ from \(i_0\). The degree and lower \(b_t\)
bounds come from the
\releasecite{Formalism/LocalRingCardinality/DataParams₂.lean}{76}{size data}.
The upper \(b_t\) bound uses the
\releasecite{Formalism/SharpRestriction/FamCore.lean}{123}{common upper endpoint}
and the
\releasecite{Formalism/SharpRestriction/ElevenLe.lean}{68}{bound on the endpoint parameter}.

Use the coordinates \(r\in P=\{(0,i,s):i\in\mathbf Z/2\mathbf Z,\ s\in H\}\)
for the stage weights. In the original \(q,d\) sign order, the CM
embedding at \((0,i,s)\) has signs \(((-1)^i,1)\) and last component
\(s\). The prescribed CM type has exactly one embedding with each
such label. Write \(\tau_r\) for the embedding with label \(r\)
at \(i(u)\). The actual weight at stage \(1\le j\le m\) is the
Harish--Chandra entry \(w_{u,j}(r)=e_j(\tau_r)\) of that record,
with \(e_j\) indexed from \(1\).
\releasecite{Formalism/LocalRingCardinality/DataParams₂.lean}{73}{The stage-weight definition}.
With the common
\(\eta_{\rm wt}>0\) of the tied records, its prescribed intervals are
\[
\begin{array}{c|c}
 \text{position }r & \text{interval for }w_{u,j}(r)\\ \hline
 (0,1,s+u),\ s\in T_2 & [b,b+\eta_{\rm wt}+1]\\
 (0,0,t),\ t\in T & [b_t+b'/x,b_t+b'/x+\eta_{\rm wt}]\\
 \text{all other positions in }P & [w'',w''+\eta_{\rm wt}].
\end{array}
\]
The first two rows are disjoint in their second coordinates.

\leavevmode\teamanchor{proof.fourier-weight-intervals}%
For the classical target the intervals follow from the
\hyperlink{proof.section-run-family-identification}{row certificate at \(i(u)\)}.
At a label \(\tau=\tau_r\) it gives the stage weight
\(e_{\tau,j-1}=B_\tau+(j-\tfrac12)H_{\rm node}\) from the
\hyperlink{proof.section-run-design}{chosen row construction}.
With \(\eta_{\rm wt}=mH_{\rm node}\) and \(H_{\rm node}\ge2\), every
stage satisfies
\[
 B_\tau+1\le w_{u,j}(r)\le B_\tau+\eta_{\rm wt}-1 .
\]
At a thick label, the start is \(\lceil b\rceil\) with a possible
torsion shift of at most \(1\), since
\hyperlink{proof.classical-host-roots-of-unity}{\(\mu(E)=\{\pm1\}\)}
on the classical host at large scales. Thus \(b\le B_\tau\le b+2\),
yielding the first interval. At a top label
we have \(B_\tau=b_t+\lfloor b'/x\rfloor\), which is between
\(b_t+b'/x-1\) and \(b_t+b'/x\). At a remaining label
\(w''\le B_\tau=\lceil w''\rceil<w''+1\). Applying the same
node inequalities gives the second and third intervals.
The classical rank ranges in Section~\ref{sec:family} give
\(n\le x^{9/4}\) and \(x^4\le k\le x^5\) for large \(x\).
The node spacing uses a fixed exponent, so that
\(\log\eta_{\rm wt}=O(\log x)\le x^{\varepsilon/4}\).
For the Dirichlet extension use the corresponding selected
\releasecite{Formalism/LocalRingCardinality/DataParams₂.lean}{151}{three weight ranges}
and the uniform
\releasecite{Formalism/SharpRestriction/ElevenLe.lean}{24}{rank ranges}
and \releasecite{Formalism/SharpRestriction/ElevenLe.lean}{81}{node allowance}.

In particular, the common size bounds give
\[
 \frac{2n\delta_1 x^\varepsilon L_{\mathrm{a},G}}{k}
 \le 2x^{-7/4+(201/100)\varepsilon}\longrightarrow0,\qquad
 w''=\frac{\exp(x^\varepsilon-x^{\varepsilon/3})}{n}\longrightarrow\infty,
 \qquad \frac{\eta_{\rm wt}}{b'}\longrightarrow0 .
\]
Thus all weights are \(\ge1\) at large \(x\). Also \(w''/b'=1/(xn)\le1/x\),
and \(b_t/b'\le1+o(1)\), so both of the last two rows have upper bound
\(2b'\) at a common threshold. Since \(a_t/w_{u,j}(r)\le2a_E\), the
cutoff \(N=\lceil2a_E\rceil\) in \eqref{eq:prime-power-cutoff} applies to
the actual degrees and weights. These same scales give
\(kb'\le a_E\) since \(k\le x^5\le X'\) for large \(x\), and
\(X''=a_E/w''\le a_E\).

Put \(\vartheta=\exp(-c_6x^{\varepsilon/8})\) with the host
constant \(c_6>0\) from \S\ref{sec:family}. Since \(f=2n\) is
polynomially bounded and \(\eta_{\rm wt}\le\exp(x^{\varepsilon/4})\),
both \(6f(x+1)(\eta_{\rm wt}+1)/b\) and \(2f(x+1)\vartheta\) are
\(o(x^{-1})\). Thus at a common threshold
\begin{equation}
\label{eq:fourier-weight-costs}
 6f(x+1)(\eta_{\rm wt}+1)\le b/x,\qquad
 2f(x+1)\vartheta\le1/x .
\end{equation}
Each threshold here follows from the fixed precision and the uniform
record bounds. In particular it applies to every \(j\) and \(u\) at
one scale, using the actual \(a_t,b_t\) from the common seed.

\subsection{Finite Fourier inversion}
\teamanchor{proof.fourier-inversion-thin}

Use the group \(\Gamma\), the subgroup \(P=\ker c\) of order \(f=2|H|\),
and the bijection \(p_t:\Gamma\setminus P\to P\) in
\eqref{eq:prime-power-ramp}. The base field is a subfield of
\(\mathbf Q(\zeta_{Q_Sqd})\), and its Artin map is the quotient of the
cyclotomic action. For \(\xi\in\widehat\Gamma\), let \(\chi_\xi\) be the
unique primitive Dirichlet character inducing \(\xi\circ\operatorname{Art}\)
on units modulo \(Q_Sqd\). Its conductor divides \(Q_Sqd\), and for
\(\xi=1\) its conductor is \(1\). Write \(G_\xi=G_{\chi_\xi}\) for the
centered sum of \eqref{eq:centred-prime-sums}, with
\eqref{eq:principal-ramp} for the trivial character. For \(\xi\ne1\)
put \(\ell_\xi=L'(1,\chi_\xi)/L(1,\chi_\xi)\), using the nonvanishing
at \(1\) in \S\ref{sec:family}.

For a one-dimensional character \(\xi\) of the Galois group of a finite
abelian extension of \(\mathbf Q\), let \(V_\xi=\mathbf C\) with the action by
\(\xi\). At each prime \(p\), let \(I_p\) be inertia and let \(F_p\)
lift the arithmetic Frobenius, which sends \(z\) to \(z^p\) on the
residue field. To compare the two Euler factors for our field, write
\(Q_Sqd=p^\alpha N_p\) with \(p\nmid N_p\).
By the cyclotomic decomposition of Section~\ref{sec:family}, inertia
is the image of \((\mathbf Z/p^\alpha\mathbf Z)^\times\) in \(\Gamma\),
and a Frobenius lift can be represented by a unit whose component
modulo \(N_p\) is \(p\bmod N_p\).
Thus \(\xi\) is trivial on \(I_p\) exactly when its pullback factors
through \(N_p\), or equivalently \(p\nmid\operatorname{cond}(\chi_\xi)\).
In that case every lift has \(\xi(F_p)=\chi_\xi(p)\).
When \(p\mid\operatorname{cond}(\chi_\xi)\), the invariant space
is zero and \(\chi_\xi(p)=0\). Both cases give
\[
 L_p(s,\xi)=\det\!\left(1-p^{-s}F_p\mid V_\xi^{I_p}\right)^{-1}
           =\left(1-\chi_\xi(p)p^{-s}\right)^{-1}
           \qquad(\Re s>1).
\]
The logarithm given by the series in \(p^{-s}\) is
\[
 \log L_p(s,\xi)=\sum_{h\ge1}\frac{v_\xi(p^h)}{h\,p^{hs}},
 \qquad
 v_\xi(p^h)=\frac{1}{|I_p|}\sum_{i\in I_p}\xi(hF_p+i)
           =\xi(F_p)^h\frac{1}{|I_p|}\sum_{i\in I_p}\xi(i)
           =\chi_\xi(p^h).
\]
The inertia average equals \(1\) on characters trivial on \(I_p\)
and vanishes on the others. This uses the primitive character even at
primes dividing \(Q_Sqd\) whose local \(p\)-part is trivial.
The coset \(hF_p+I_p\) is independent of the chosen lift because
changing \(F_p\) by an element of \(I_p\) changes \(hF_p\) by an
element of \(I_p\).
\cite[(3.2)--(3.5), (3.15), and \S4.2]{deshalit-artin}

For a finite abelian group \(B\) written additively, the dual group
\(\widehat B\) has \(|B|\) characters. For every \(z\in B\) and
\(\xi\in\widehat B\), the two orthogonality relations are
\[
 \frac{1}{|B|}\sum_{\xi'\in\widehat B}\xi'(z)=\mathbf 1_{\{z=0\}},
 \qquad
 \frac{1}{|B|}\sum_{z'\in B}\xi(z')=\mathbf 1_{\{\xi=1\}}.
\]
On each cyclic factor the first sum is a geometric series. For the
second, translating the sum by an element with character value
different from \(1\) forces it to vanish for a nontrivial character.
Consequently, for every \(F:\Gamma\to\mathbf C\) and \(g\in\Gamma\),
Fourier inversion in the convention used here is
\[
 F(g)=\frac{1}{2f}\sum_{\xi\in\widehat\Gamma}\overline{\xi(g)}
                                \sum_{z\in\Gamma}F(z)\xi(z).
\]
\cite[Theorem~4.4, (4.12)--(4.13)]{z6montgomery-vaughan}
Apply this with \(F(z)=|I_p|^{-1}\mathbf 1_{\{z-hF_p\in I_p\}}\).
The inner sum is \(v_\xi(p^h)\), giving
\[
 \frac{1}{f}\sum_{\xi\in\widehat\Gamma}\overline{\xi(g)}v_\xi(p^h)
 =\frac{2}{|I_p|}\mathbf 1_{\{g-hF_p\in I_p\}} .
\]
Multiplying by \(\Lambda(p^h)=\log p\) on \(\Gamma\setminus P\)
recovers the coefficient in \eqref{eq:prime-power-ramp}. When
\(p\mid Q_Sq\), both the Frobenius lift and inertia lie in \(P\) since
\(p\) splits in \(\mathbf Q(\sqrt{-d})\). For the classical host these
subgroups are explicit in \eqref{eq:classical-host-inertia}. At \(p=d\),
the two elements of \(hF_d+I_d\) have opposite \(c\)-values, so there
is one element outside \(P\) for each \(h\ge1\), of weight \(\log d\).
For \(p\nmid Q_Sqd\), the subgroup \(I_p\) is trivial. The unique
term of the coset has weight \(2\log p\) if \(c(F_p)=-1\) and \(h\) is
odd. These are the support conditions used in Section~\ref{sec:integral}.

For \(r=p_t(g)\), factor the positive part as
\((a/p^h-w_r)_+=w_r((a/w_r)/p^h-1)_+\). At \(a=a_t\) and
\(w=w_{u,j}\), the common cutoff \(N\) in \eqref{eq:prime-power-cutoff}
satisfies \(a/w_r\le N\). Thus the sum through \(N\) equals
the unrestricted sum in \eqref{eq:prime-power-ramp} used to define \(D(u)\).
The weighted uncentered sum is
\(w_rG_\xi(a/w_r)+a(\log a-\log w_r-1-\gamma)\) for \(\xi=1\),
and \(w_rG_\xi(a/w_r)-a\ell_\xi\) for \(\xi\ne1\).
Writing \(\Gamma\setminus P=(1,0,0)+P\), the mean of
\(\overline{\xi(g)}\) on this coset is \(\overline{\xi(1,0,0)}\)
times its mean on \(P\). It vanishes unless \(\xi|_P=1\).
The two remaining characters are \(1,c\), with means \(1,-1\)
on \(\Gamma\setminus P\). The bijection \(p_t\) identifies the sum of \(\log
w_{p_t(g)}\) over \(\Gamma\setminus P\) with the sum of \(\log w_r\)
over \(P\). The uncentered sum therefore equals the sum of weighted
\(G_\xi\)'s plus
\[
 a(\log a-1-\gamma)-\frac{a}{f}\sum_{r\in P}\log w_r
       -a(-1)\ell_c=M(a,w).
\]
Here the last term uses the same primitive \(\eta_d\) as the CM
normalization in \eqref{eq:prime-power-center}. We get
\begin{equation}
\label{eq:fourier-balance}
 A(a,t,w)-M(a,w)=\frac{1}{f}\sum_{g\notin P}
   \sum_{\xi\in\widehat\Gamma}\overline{\xi(g)}\,w_{p_t(g)}
                      G_\xi\bigl(a/w_{p_t(g)}\bigr).
\end{equation}
The left side is real since \(A(a,t,w)\) is real and
\(\ell_c=L'(1,\eta_d)/L(1,\eta_d)\) is real. In particular,
the real part of the right side at \(a=a_t\), \(w=w_{u,j}\),
summed over \((t,j)\), is exactly \(D(u)\) in
\eqref{eq:centered-excess-definition}.
\releasecite{Development/AutomorphicArithmetic/DedekindLFunctions/LandauTheorems/ChBalance.lean}{200}{Formal identity}.

\subsection{Reduction to the large weights}

Let \(H_0\le H\) be the distinguished cyclic subgroup from
\S\ref{sec:family}, with \(o=[H:H_0]\) equal to the order of the target
character \(\psi\). Embed \(H_0\) into \(\Gamma\) by \(h\mapsto(0,0,h)\).
Put \(\delta_2=N_2/n\). For each \(u\in H_0\) and stage \(j\),
the positions with \(s\in T_2\subset H_0\) in the first row of the
weight table are distinct. Section~\ref{sec:fourier-sampling}
places their weights in \([b,b+\eta_{\rm wt}+1]\) and all other
weights in \([1,2b']\).

For every real \(z>0\), the Chebyshev function satisfies
\(\mathcal V(z):=\sum_{n\le z}\Lambda(n)<1.03883\,z<3z\).
\cite[Theorem~12, (3.35)]{rosser-schoenfeld}
Write \(z=a_t/w\). In the nonprincipal case
\(wG_\xi(a_t/w)=\sum_{n\ge1}\Lambda(n)\chi_\xi(n)(a_t/n-w)_++a_t\ell_\xi\).
The last term is constant in \(w\). Between ramp breakpoints the
derivative is \(-\sum_{n<z}\Lambda(n)\chi_\xi(n)\) for \(\xi\ne1\), and
\(z-\sum_{n<z}\Lambda(n)\) for \(\xi=1\). In both cases its modulus is
at most \(z+\mathcal V(z)<3z\). On the replacement interval \(w\ge b\)
we have \(z\le x+1\). Continuity at the breakpoints bounds the change
by \(3(x+1)(\eta_{\rm wt}+1)\).

\leavevmode\teamanchor{proof.thin-position-bound}%
For a small weight \(w\), \(b'X'=xb\) and \(x\ge3\) imply
\(a_t/w\ge (x-1)b/(2b')\ge X'/3\). The base host satisfies
\[
 |G_\xi(y)|\le \vartheta y\qquad
 (\xi\in\widehat\Gamma,\ y\ge X'/3).
\]
Thus each omitted term has modulus at most \((x+1)b\vartheta\).
The estimate covers every \(\xi\), because the small weights are
omitted before taking the average over \(H_0\). The sum in
\eqref{eq:fourier-balance} has \(f\) positions, \(2f\) characters
and prefactor \(1/f\). Multiplying each of the two error bounds by
\(2f\) and applying \eqref{eq:fourier-weight-costs} gives a total
error at most \(2b/x\) for each \((t,j,u)\). This uses the tail
\eqref{eq:host-character-tail} at \(i_0\).

\subsection{The translate average}
\teamanchor{proof.translate-average-zeta}

Write \(e(h)=(1,0,h)\). The preimage of \((0,1,s+u)\) under \(p_t\)
is \(e(s+u-t)\). Its conjugate character value is
\[
 \overline{\xi(e(s+u-t))}
   =\xi(e(t))\,\overline{\xi(0,0,s)}\,\overline{\xi(0,0,u)}.
\]
Indeed \(e(s+u-t)=-e(t)+(0,0,s)+(0,0,u)\) since the first coordinate
has order two.
By character orthogonality on \(H_0\), the mean over \(u\) vanishes unless
\(\xi\in H_0^\perp=\{\xi:\xi|_{H_0}=1\}\). For such a character,
\(s\in H_0\) also gives a factor of \(1\), so all \(N_2\) large
weights contribute equally. Put \(Q=N_Tmb\) and
\(\overline D_x=|H_0|^{-1}\sum_{u\in H_0}D(u)\).
The replacement estimate bounds the real part of the reduced sum
at a fixed \(u\) by \(D(u)+2Q/x\), since there are \(mN_T\) pairs
\((t,j)\). Each replaced term uses \(bG_\xi(y_t)\),
which is independent of the stage and the translate. Averaging the
real inequality gives
\begin{equation}
\label{eq:fourier-translate-mean}
 \frac{Q\delta_2}{2}\operatorname{Re}
   \frac{1}{N_T}\sum_{t\in T}\sum_{\xi\in H_0^\perp}
                 \xi(e(t))G_\xi(y_t)
 \le \overline D_x+\frac{2Q}{x}.
\end{equation}
The factor \(Q\delta_2/2=mN_TbN_2/f\) comes from the \(m\) stages,
the \(N_2\) positions and the \(1/f\) in \eqref{eq:fourier-balance}.
It is positive since \(N_T,m,b>0\) and \(\delta_2\ge x^{-\varepsilon}\).
Here \(\overline D_x\) is the uniform mean of the signed centered
statistic on all the tied representatives at this one real scale.
For classical zeta \(H_0=H\) and \(\overline D_x=\langle D\rangle\).
The lattice bound
\eqref{eq:classical-averaged-excess} therefore bounds the right-hand
side by \(E_x+2Q/x\). For the Dirichlet extension the same upper
bound follows by averaging \eqref{eq:integral-excess} over \(H_0\).

\rhrange{The release's application of the Fourier theorem uses a
separate upper bound at each translate. In the matched normalization
its inequality for \texttt{gapA} reads \(D(u)\le E_x\) for every
translate at the fixed classical scale.
\releasecite{Development/AutomorphicArithmetic/DedekindLFunctions/SplitInertia.lean}{50}{The centered quantity}
\releasecite{Development/AutomorphicArithmetic/DedekindLFunctions/LandauTheorems/ChTop.lean}{37}{The pointwise bound}
\releasecite{Development/AutomorphicArithmetic/PrimitiveCosets/CoreGm.lean}{55}{Its actual Fourier application}}
{The single inequality \(\overline D_x=\langle D\rangle\le E_x\)
obtained in \eqref{eq:classical-averaged-excess} by summing the
classical joint-deficit inequality. Summing the Fourier comparisons
retains the signed mean on the right of
\eqref{eq:fourier-translate-mean}. Orthogonality evaluates the left.}
The release also has a conditional estimate of the whole-group
mean of the same statistic, \texttt{topExcess}, under its own
slope and loss hypotheses.
\releasecite{Formalism/SizeWindowBrackets/GrowthBoundGivesRiemannHypothesis.lean}{50}{The whole-group mean}
\releasecite{Formalism/SizeWindowBrackets/ChainSlopeWindows/PChain.lean}{45}{Its bound under those hypotheses}.

\subsection{Selecting the target characters}

Restriction identifies \(H/H_0\) with the Galois group of the
cyclotomic subfield \(\mathbf Q(\psi)\) fixed by \(\ker\psi\).
Its Artin map is the quotient by \(\ker\psi\) of
\((\mathbf Z/M\mathbf Z)^\times\). Since \(\psi\) is faithful on that
cyclic quotient, the characters of \(H/H_0\) pull back along
\(h_{\rm nat}\) to the powers \(\psi^k\), \(0\le k<o\).
View each such character \(\nu\) on \(H\). As \(\nu\) and
\(i,j\in\{0,1\}\) vary, the factor
\(\nu(h_-)^{e_q}\) permutes the sign choices since \(2h_-=0\). Thus all
\(4o\) characters in \(H_0^\perp\) have the form
\[
 \xi(e_d,e_q,h)=(-1)^{i e_q+j e_d}\nu(h)\nu(h_-)^{e_q}.
\]
Substituting \eqref{eq:artin-coordinates}, the pullback at a unit
\(a\bmod Q_Sqd\) is
\[
 \eta_q(a)^i\eta_d(a)^j\,\nu(h_{\rm nat}(a))\,\nu(h_-)^{2e_q(a)}
 =\psi^k(a)\eta_q(a)^i\eta_d(a)^j .
\]
Its primitive associate is \(\chi_k\eta_q^i\eta_d^j\), where \(\chi_k\)
induces \(\psi^k\) modulo \(M\) and is primitive, with \(\chi_0=1\).
Omit a quadratic factor with exponent zero. For a product of primitive
characters of pairwise coprime conductors, the Chinese remainder theorem
lets us vary the unit in one factor while setting every other component
to \(1\). If the product factors through a divisor of the product
modulus, each original factor must then factor through the part of this
divisor supported on its own conductor. By primitivity that part equals
its full conductor. Since \(q,d\) are coprime to \(M\), the conductor here
is \(\operatorname{cond}(\chi_k)q^id^j\). Thus every term
with \((i,j)\ne(0,0)\) satisfies the nontrivial quadratic twist hypotheses
of Section~\ref{sec:family}.
At \(e(t)=(1,0,t)\) the phase is \((-1)^j\nu(t)\).
Let \(\mathcal X\) consist of the characters with \(i=j=0\).
Their inducing primitive characters form
the fixed set \(\mathcal K\) attached to the
\releasecite{Formalism/SharpRestriction/RamificationFrobenius/TotientGaloisClasses/ExplicitFormulaErrorOfHostConditions.lean}{105}{subgroup generated by \(\psi\)}.
Different powers give different primitive associates since an associate
determines its inducing values on units modulo \(M\). Thus \(\mathcal K\)
has \(o\) elements, contains \(1,\psi\) and is closed under conjugation.
The phase \(\xi(e(t))\) is \(\nu(t)\) for \(\xi\in\mathcal X\), and
\(\nu\) is nontrivial when \(\xi\ne1\).

For \(v\in H/H_0\), the proportion of \(T\) in \(v\) differs from
\((3+\mathbf 1_{\{v=0\}})/(3o+1)\) by at most \(x^{-1+\varepsilon}\).
For a nontrivial character \(\nu\) of the quotient, orthogonality gives
\[
 \frac{3\sum_{v\in H/H_0}\nu(v)+\nu(0)}{3o+1}=\frac{1}{3o+1}.
\]
Thus \(N_T^{-1}\sum_{t\in T}\xi(e(t))\) differs from
\((3o+1)^{-1}\) by at most \(o x^{-1+\varepsilon}\) for a nontrivial
retained character. For \(\xi=1\) the coefficient is exactly \(1\).
Put \(a_1=1\) and \(a_\chi=(3o+1)^{-1}\) for
\(\chi\in\mathcal K\setminus\{1\}\), and set
\[
 G_{\mathcal K}(x)=\sum_{\chi\in\mathcal K}a_\chi\Re G_\chi(x).
\]

For every real \(z\ge2\), Mertens' estimate gives
\(\sum_{n\le z}\Lambda(n)/n=\log z+O(1)\), with an absolute implied
constant.
\cite[Theorem~2.7(a)]{z6montgomery-vaughan}
Indeed, for \(v=\lfloor z\rfloor\), summing
\(\log n=\sum_{d\mid n}\Lambda(d)\) up to \(v\) gives
\(\log(v!)=\sum_{d\le v}\Lambda(d)\lfloor v/d\rfloor\).
Replacing each floor by \(v/d\) costs at most \(\mathcal V(v)\ll v\).
The integral comparison for \(\log(v!)\) gives \(\log(v!)/v=\log z+O(1)\).
Between breakpoints, the derivative of the uncentered ramp sum at \(y\)
is \(\sum_{n<y}\Lambda(n)\chi_\xi(n)/n\), of modulus at most
\(\log y+O(1)\) since \(|\chi_\xi(n)|\le1\). For \(\xi=1\) the
centering adds \(-\log y+\gamma\) to the derivative. For \(\xi\ne1\)
it adds \(\ell_\xi\), which is fixed for
\(\mathcal K\setminus\{1\}\) and is \(O(\log x)\) for the auxiliary
twists by \S\ref{sec:family}. For each \(\xi\in H_0^\perp\) these
continuous centered sums therefore satisfy
\(|G_\xi(y_t)-G_\xi(x)|=O(\log x)\) on \([x-1,x+1]\).
The same Mertens estimate gives an uncentered sum of modulus
at most \(x(\log x+O(1))\), so \(|G_\xi(x)|=O(x\log x)\). Replacing
\(y_t\) by \(x\) in the normalized sum costs \(O(o\log x)\). Replacing the
mean phase of a retained nontrivial character by its limit costs
at most \(o x^{-1+\varepsilon}|G_\xi(x)|\).
The combined error is \(O(x^\varepsilon\log x)\) for fixed target
and precision, hence at most \(4x^{1/2+\varepsilon}\) for large \(x\).
Write
\(\Sigma(x)=\sum_{\xi\in H_0^\perp\setminus\mathcal X}|G_\xi(x)|\).
Each term with \(\xi\notin\mathcal X\) has real part at least
\(-|G_\xi(x)|\), since its multiplier is the mean of unit complex
numbers. The real character sum averaged over \(t\) on the left of
\eqref{eq:fourier-translate-mean}, before multiplication by
\(Q\delta_2/2\), is therefore at least
\(G_{\mathcal K}(x)-\Sigma(x)-4x^{1/2+\varepsilon}\). This gives
\begin{equation}
\label{eq:fourier-extraction}
 \frac{Q\delta_2}{2}
       \bigl(G_{\mathcal K}(x)-\Sigma(x)-4x^{1/2+\varepsilon}\bigr)
 \le \overline D_x+\frac{2Q}{x}.
\end{equation}
The bound \(|G_\xi(x)|\le x^{1/2+\varepsilon}\) for these twists
gives \(\Sigma(x)\le4o\,x^{1/2+\varepsilon}\).

\leavevmode\teamanchor{proof.zeta-specialization}%
For \(\psi=1\), \(o=1\) and \(H_0=H\). The expression in the
translate average at each \(t\) has the four terms
\[
 G_1(y_t)+G_{\eta_q}(y_t)-G_{\eta_d}(y_t)-G_{\eta_q\eta_d}(y_t).
\]
The signs are the evaluations at \((1,0,0)\). The three auxiliary
characters have conductors \(q,d,qd\) and enter \(\Sigma(x)\) through
the quadratic twist bounds. In particular the classical average is
at least \(G_1(x)-3x^{1/2+\varepsilon}-C\log x\), with a constant
uniform in the selected \(q,d\), since \(|\ell_\xi|\le4\log x\) for
those twists. This calculation uses the size of \(T\) and the
degree interval. The retained sum is \(G_{\mathcal K}=G_1\).

\subsection{The square-root scale}
\teamanchor{proof.fourier-exponent}

The lattice bound just substituted also bounds the right-hand side
of \eqref{eq:fourier-extraction} by \(E_x+2Q/x\).
Write \(C_{\rm err}=e_D+e_L+3\). For \(\psi=1\) the budget is
\(E_x=B^+x^{5\varepsilon}\) from \eqref{eq:classical-averaged-excess}.
For the Dirichlet extension it is the budget of \eqref{eq:integral-excess}.
Dividing by \(Qx^{C_{\rm err}\varepsilon}\) gives
\[
 \frac{a_E}{mb}+\frac{a_E m}{bk}
     +2\frac{n}{N_T}m x^\varepsilon
     +\frac{N_2}{N_T}\frac{k}{n^2}x^{2\varepsilon}.
\]
In addition to the common bounds \eqref{eq:common-error-window-sizes},
we have \(a_E=xb\), \(N_T\ge nx^{-\varepsilon/8}/2\), and
\[
 x^{1/2-\varepsilon}\le m,\qquad k\le n^2x^{1/2-\varepsilon}.
\]
For \(\psi=1\), the sample counts are \eqref{eq:classical-sampling-size}
and the bounds for \(m,k\) follow from the prescribed ranks in
Section~\ref{sec:family}. The same size inequalities hold for the
\releasecite{Formalism/SharpRestriction/GlobalSchurBounds/LossesSmallShares.lean}{281}{selected Dirichlet data}.
Also \(N_2\le n\) because \(T_2\subseteq H\).
The first and third terms have bounds \(2x/m\) and
\(4mx^{9\varepsilon/8}\), which balance at the square-root scale
up to powers \(x^{O(\varepsilon)}\). The four bounds in order are
\(2x^{1/2+\varepsilon}\), \(2\),
\(4x^{1/2+13\varepsilon/8}\), and \(2x^{1/2+9\varepsilon/8}\).
Their sum is at most \(10x^{1/2+13\varepsilon/8}\).
Divide the resulting estimate by \(Q\delta_2/2>0\) and use
\(\delta_2\ge x^{-\varepsilon}\).
The factors from the main budget have exponent loss
\(C_{\rm err}+13/8+1=C_{\rm err}+21/8\), giving
\[
 G_{\mathcal K}(x)-\Sigma(x)\le
 20x^{1/2+(C_{\rm err}+21/8)\varepsilon}
 +4x^{-1+\varepsilon}+4x^{1/2+\varepsilon}
 \le 26x^{1/2+(C_{\rm err}+21/8)\varepsilon}.
\]
In the last bound use \(4x^{-1+\varepsilon}\le2x^{1/2+\varepsilon}\)
at \(x\ge3\), and \(C_{\rm err}\ge0\).
For sufficiently large \(x\), \(26\le x^{3\varepsilon/8}\) and
\(1+4o\le x^\varepsilon\). The first inequality bounds
\(G_{\mathcal K}(x)-\Sigma(x)\) by \(x^{1/2+(C_{\rm err}+3)\varepsilon}\).
The bound on \(\Sigma(x)\) brings the coefficient to at most \(1+4o\).
Absorbing it gives
\begin{equation}
\label{eq:fourier-pointwise}
 G_{\mathcal K}(x)\le x^{1/2+(e_D+e_L+7)\varepsilon} .
\end{equation}
The threshold in this calculation depends on the target, fixed
precision and uniform comparison constants. For \(\psi=1\),
\(\mathcal K=\{1\}\) with weight \(1\) and \((e_D,e_L)=(2,0)\).
With the signed mean from Section~\ref{sec:integral} and the classical
sampling and auxiliary character estimates, \eqref{eq:fourier-pointwise}
gives \(G_1(x)\le x^{1/2+9\varepsilon}\) at this same \(x\).
Section~\ref{sec:assembly} applies this conclusion throughout the
real intervals of the fixed family.
